\chapter{Marco tecnológico y estado del arte}

Este capítulo sitúa el trabajo en su contexto técnico: las tecnologías sobre
las que se construye el framework y los enfoques existentes para el problema
de generar interfaces a partir de un modelo de datos.

\section{Tecnologías del backend}

\subsection{FastAPI y Pydantic}
FastAPI~\cite{fastapi} es un \emph{framework} web asíncrono para Python que
genera documentación OpenAPI automáticamente y valida las entradas y salidas
mediante Pydantic. La elección es relevante para este trabajo porque Pydantic
puede emitir \gls{jsonschema}~\cite{jsonschema} a partir de sus modelos, lo
que convierte a FastAPI en un punto natural para publicar el esquema de cada
entidad. La inyección de dependencias de FastAPI se emplea para gestionar las
sesiones de base de datos y la autenticación de forma transversal.

\subsection{SQLAlchemy 2.x y Alembic}
SQLAlchemy~\cite{sqlalchemy} es el ORM de referencia en Python. Su versión
2.x introduce la API tipada con \texttt{Mapped[...]} y \texttt{mapped\_column},
que el framework aprovecha para inferir tipos al construir el esquema. Alembic
gestiona las migraciones de esquema, imprescindible cuando el modelo de datos
evoluciona: cada cambio se materializa en una \emph{revision} versionada y
reproducible.

\subsection{SQLAlchemy-Continuum}
\gls{continuum} añade versionado e histórico automático a los modelos ORM. Su
particularidad --- debe inicializarse antes de definir los modelos --- es una
restricción de diseño que condicionó la arquitectura del inicializador. El
versionado nativo delega el histórico en el propio motor cuando es posible,
reduciendo el coste de mantener tablas espejo.

\subsection{Infraestructura: PostgreSQL, Docker, Redis y Celery}
PostgreSQL es la base de datos relacional elegida por su robustez y su soporte
de tipos avanzados (JSONB, \emph{full-text search}). \gls{docker} y
\texttt{docker compose} aportan despliegue reproducible multiplataforma
(RNF-02). Redis y Celery habilitan el procesamiento asíncrono de tareas
(\texttt{SystemProcess}), inicializados de forma opcional por el framework.

\section{Tecnologías del frontend}

\subsection{Angular}
Angular~\cite{angular} es un \emph{framework} de aplicaciones web basado en
TypeScript, con inyección de dependencias y un sistema de módulos que encaja
con la distribución del proyecto como librería publicable (\texttt{ngt-gui}).
Su tipado fuerte permite que la capa cliente del contrato sea genérica y
tipada.

\subsection{Formly}
\gls{formly}~\cite{formly} construye formularios de manera declarativa a
partir de una configuración (\texttt{FormlyFieldConfig[]}) en lugar de HTML
campo a campo. Soporta tipos personalizados, validadores y expresiones
reactivas, lo que lo hace idóneo como destino de una traducción automática
desde el esquema. El proyecto lo combina con \emph{ng-zorro} y un conjunto de
tipos personalizados.

\section{Estado del arte: generación de interfaces dirigida por esquema}
La idea de derivar interfaces de un esquema no es nueva. En el ecosistema web
existen \texttt{react-jsonschema-form} y \texttt{JSON Forms}, que renderizan
formularios a partir de \gls{jsonschema}. Estas soluciones, sin embargo, se
centran en el lado cliente y asumen que el esquema se proporciona ya
elaborado. En el lado servidor, herramientas como FastAPI exponen OpenAPI,
pero ese contrato describe \emph{operaciones}, no incluye pistas de
presentación, y no resuelve el acoplamiento del cliente al dominio.

La tabla~\ref{tab:soa} compara los enfoques. La aportación diferencial de este
trabajo no es inventar un renderizador de esquemas, sino \textbf{cerrar la
cadena completa}: derivar el esquema \emph{desde el ORM}, enriquecerlo con
pistas de presentación en el propio esquema, transportarlo por un contrato
uniforme con versionado, y consumirlo con una capa genérica que mantiene la
compatibilidad con el código heredado.

\begin{table}[!htbp]
\caption{Comparación de enfoques de generación de interfaces.}
\label{tab:soa}
\centering
\begin{tabular}{|p{3.4cm}|c|c|c|c|}
\hline
\textbf{Enfoque} & \textbf{Desde ORM} & \textbf{Hints UI} &
\textbf{Contrato} & \textbf{Versionado} \\ \hline
react-jsonschema-form & No & Parcial & No & No \\ \hline
JSON Forms & No & Sí (uischema) & No & No \\ \hline
OpenAPI/Swagger UI & No & No & Operaciones & No \\ \hline
Admin autogenerados & Sí & Limitado & Acoplado & No \\ \hline
\textbf{Este trabajo} & \textbf{Sí} & \textbf{Sí (x-*)} &
\textbf{Uniforme} & \textbf{ETag} \\ \hline
\end{tabular}
\end{table}

Frente a los \emph{admin} autogenerados (al estilo del de Django), que sí
parten del ORM pero acoplan la interfaz al servidor y son difíciles de
reutilizar fuera de su \emph{stack}, la solución propuesta mantiene frontend y
backend desacoplados mediante un contrato explícito y portable, que es
precisamente lo que habilita la reutilización en proyectos con modelos de
datos en evolución.
